% ==============================================================================
% 00_config.tex - PREÂMBULO E CONFIGURAÇÕES GLOBAIS
% Artigo: A Contribuição de Federico Faggin para a Revolução dos Microprocessadores
%         e seus Impactos Sociais: Uma Análise Histórico-Pedagógica
% Padrão: IEEE Transactions (Two-Column Journal Format)
% ==============================================================================

% ------------------------------------------------------------------------------
% 1. CODIFICAÇÃO E SUPORTE A IDIOMA
% ------------------------------------------------------------------------------
\usepackage[utf8]{inputenc}               % Suporte à codificação UTF-8
\usepackage[T1]{fontenc}                  % Codificação de fontes de 8 bits
\usepackage[brazilian]{babel}             % Hifenização e termos em Português do Brasil

% Adequação de termos pré-definidos do IEEEtran para o Português
\renewcommand\IEEEkeywordsname{Palavras-Chave}
\addto\captionsbrazilian{%
  \renewcommand{\figurename}{Figura}%
  \renewcommand{\tablename}{Tabela}%
  \renewcommand{\refname}{Referências}%
  \renewcommand{\abstractname}{Resumo}%
}

% ------------------------------------------------------------------------------
% 2. PACOTES MATEMÁTICOS E FÍSICOS
% ------------------------------------------------------------------------------
\usepackage{amsmath, amssymb, amsfonts}  % Símbolos e ambientes matemáticos avançados
\usepackage{siunitx}                     % Unidades do SI e notação física padronizada
\sisetup{
  output-decimal-marker = {,},
  exponent-product = \cdot,
  per-mode = symbol
}

% ------------------------------------------------------------------------------
% 3. ELEMENTOS GRÁFICOS E TABELAS
% ------------------------------------------------------------------------------
\usepackage{graphicx}                    % Inclusão de figuras e diagramas
\graphicspath{{img/}}                    % Diretório padrão de imagens

\usepackage{booktabs}                    % Linhas tipográficas profissionais (\toprule, \midrule, \bottomrule)
\usepackage{tabularx}                    % Tabelas com largura controlada e quebra automática de parágrafo
\usepackage{array}                       % Formatação customizada de colunas
\usepackage{makecell}                    % Quebras de linha dentro de células
\usepackage{multirow}                    % Células de múltiplas linhas

% ------------------------------------------------------------------------------
% 4. CITAÇÕES, REFERÊNCIAS E TIPOGRAFIA
% ------------------------------------------------------------------------------
\usepackage{cite}                        % Compressão e ordenação de citações [1]-[3]
\usepackage{url}                         % Formatação de URLs
\usepackage{microtype}                   % Ajustes microtipográficos (evita overfull \hbox)

% Hiperlinks clicáveis sem caixas delimitadoras coloridas (estética clássica IEEE)
\usepackage[hidelinks]{hyperref}

% Referências cruzadas semânticas (cleveref) compatíveis com padrão brasileiro
\usepackage[nameinlink,noabbrev,brazilian]{cleveref}
\crefname{figure}{Figura}{Figuras}
\Crefname{figure}{Figura}{Figuras}
\crefname{table}{Tabela}{Tabelas}
\Crefname{table}{Tabela}{Tabelas}
\crefname{equation}{Equação}{Equações}
\Crefname{equation}{Equação}{Equações}
\crefname{section}{Seção}{Seções}
\Crefname{section}{Seção}{Seções}
\crefname{subsection}{Subseção}{Subseções}
\Crefname{subsection}{Subseção}{Subseções}

% ------------------------------------------------------------------------------
% 5. PARÂMETROS INSTITUCIONAIS E METADADOS
% ------------------------------------------------------------------------------
\def\universidade{Universidade Federal de Minas Gerais}
\def\siglauni{UFMG}
\def\faculdade{Escola de Engenharia}
\def\departamento{Departamento de Engenharia Elétrica}
\def\disciplina{Descobertas de Engenharia Elétrica que Moldaram a Sociedade Contemporânea}
\def\semestre{2º/2026}
\def\autorartigo{Luis Henrique}
\def\emailautor{luis@ufmg.br}

% Definição do Título e Cabeçalho do Artigo
\title{A Contribuição de Federico Faggin para a Revolução dos Microprocessadores e seus Impactos Sociais: Uma Análise Histórico-Pedagógica}

\author{\autorartigo%
\thanks{Este trabalho faz parte da disciplina: \disciplina, oferecida no \semestre\ pelo \departamento\ da \universidade\ (\siglauni).}%
\thanks{O autor é discente do curso de Engenharia Elétrica da \universidade\ (\siglauni), Av. Antônio Carlos, 6627, Pampulha, 31270-901, Belo Horizonte, MG, Brasil (e-mail: \emailautor).}%
}

% Cabeçalho de página IEEE (Running Header)
\markboth{UFMG -- DEE -- DESCOBERTAS DE ENGENHARIA ELÉTRICA}{}
